\section{Conclusion}
\label{sec:conclusion}

In this work, we study subjective visual intention recognition under supervisory ambiguity and semantic ambiguity, and propose FDIL as a functionally decoupled framework. In the current formulation, fixed semantic priors provide structured candidate semantics, a residual student supports image-aware intent prediction on top of the semantic prior, and uncertainty-aware text teachers stabilize ambiguous supervision.

Extensive experiments on Intentonomy show that the proposed framework consistently improves over strong visual baselines and prior state-of-the-art methods. \revised{Cross-dataset studies on the Flickr-LDL and Emotion6 label-distribution benchmarks show that the division of labour between the two modules persists, with the semantic prior moving pointwise distribution fit and the uncertainty-guided teacher moving rank correlation.} Overall, the results suggest that once visual representations are sufficiently strong, the key challenge of intention understanding lies in resolving supervisory and semantic ambiguity rather than in extracting low-level visual features alone. We hope this functionally decoupled formulation provides a clearer perspective for future research on human-centered visual understanding tasks.
